% SPDX-License-Identifier: Apache-2.0
\section{What a Unit States About Itself}
\label{sec:x-formal}

Every other example here is checked by comparison: the netlist must
do what the run did. That catches a lowering that is wrong, and it says
nothing about a design whose Rust is wrong in the same way as its
netlist. A unit can also say what must hold of it, where it is
written, and then a simulator checks the statement, and a formal tool
can prove it for every input rather than for the one run.

\code{check!(c, "message")} says that \code{c} holds on every edge
at which the statement is reached. \code{assume!(c, "message")} says
what the unit takes for granted of its inputs, which a formal tool
proves the checks under. \code{cover!(c, "message")} says that
\code{c} can happen. In the run a failed check stops it with its
message, and so does an assumption its own inputs break; a cover
point is counted, and \code{covered} reports how often it was reached.

In the netlist each is a statement of the clocked process it is in,
under the same conditions, after the reset. The VHDL writes a check
and an assumption as VHDL's own \code{assert}, which nvc checks while
it replays the run; the Verilog writes all three as SystemVerilog's
immediate assertions between \code{`ifdef FORMAL} and \code{`endif},
which Verilator checks here and a synthesis tool never sees. Both
co-simulations were checked to fail when a condition in the netlist
does not hold.

The counter below counts in decimal by a step of nought, one or two.
It checks that its count is a digit and, under the condition that it
wraps, that the wrap lands below two; it assumes its step is never
three, which it does not handle; and it covers the count reaching
nine.

\begin{figure}[htbp]
\centering
\input{formal_timing.tex}
\caption{The digit counting by its step and wrapping past nine; every
statement holds on every edge, and nine is reached twice.}
\label{fig:x-formal}
\end{figure}

\lstinputlisting[caption={\code{ex\_formal.rs}. Run with \code{bazel run //lib/examples:ex\_formal}.},label={lst:x-formal}]{lib/examples/ex_formal.rs}

\lstinputlisting[style=ascii,caption={What \code{ex\_formal} printed when the build ran it: the counter, the cover point's count, then the Verilog with the four statements.},label={lst:x-formal-out}]{out_formal.txt}
